% Standalone problem document, generated from the adjacent README.md.
% Compile with XeLaTeX. Mathematical content is maintained in Markdown.
\documentclass[11pt,a4paper]{article}
\usepackage[fontsize=10.5pt]{scrextend}
\usepackage[margin=22mm,headheight=15pt,headsep=9mm,footskip=11mm]{geometry}
\usepackage{fontspec}
\setmainfont{texgyrepagella}[Extension=.otf,UprightFont=*-regular,BoldFont=*-bold,ItalicFont=*-italic,BoldItalicFont=*-bolditalic]
\setsansfont{texgyreheros}[Extension=.otf,UprightFont=*-regular,BoldFont=*-bold,ItalicFont=*-italic,BoldItalicFont=*-bolditalic]
\setmonofont{texgyrecursor}[Extension=.otf,UprightFont=*-regular,BoldFont=*-bold,ItalicFont=*-italic,BoldItalicFont=*-bolditalic,Scale=MatchLowercase]
\usepackage{amsmath,amssymb}
\usepackage{unicode-math}
\setmathfont{texgyrepagella-math.otf}
\usepackage{xcolor,microtype,enumitem,needspace,fancyhdr}
\usepackage{bookmark}
\definecolor{ink}{HTML}{17344B}
\definecolor{accent}{HTML}{197D80}
\definecolor{muted}{HTML}{596773}
\hypersetup{colorlinks=true,linkcolor=accent,urlcolor=accent,pdftitle={PF-03 - Rational
factors for rational completely positive boundary
matrices},pdfauthor={Open Problems in Numerical Linear Algebra}}
\urlstyle{same}
\setlength{\parindent}{0pt}
\setlength{\parskip}{5pt plus 1pt minus 1pt}
\linespread{1.04}
\setlength{\emergencystretch}{3em}
\widowpenalty=10000
\clubpenalty=10000
\displaywidowpenalty=10000
\allowdisplaybreaks[1]
\setlist{nosep,leftmargin=1.5em}
\providecommand{\tightlist}{\setlength{\itemsep}{0pt}\setlength{\parskip}{0pt}}
\providecommand{\pandocbounded}[1]{#1}
\pagestyle{fancy}
\fancyhf{}
\fancyhead[L]{\sffamily\footnotesize\color{muted}OPEN PROBLEMS IN NUMERICAL LINEAR ALGEBRA}
\fancyhead[R]{\sffamily\bfseries\footnotesize\color{accent}PF-03}
\fancyfoot[L]{\parbox[b]{0.9\textwidth}{\sffamily\fontsize{8}{10}\selectfont\color{muted}Editorial ratings; a bounded search cannot certify that a problem remains unsolved.\\Verification check: 2026-09-19}}
\fancyfoot[R]{\sffamily\footnotesize\color{muted}\thepage}
\renewcommand{\headrulewidth}{0.3pt}
\renewcommand{\headrule}{\hbox to\headwidth{\color{accent}\leaders\hrule height \headrulewidth\hfill}}
\makeatletter
\renewcommand{\section}{\Needspace{5\baselineskip}\@startsection{section}{1}{0pt}{12pt plus 2pt minus 2pt}{4pt}{\sffamily\bfseries\large\color{ink}}}
\renewcommand{\subsection}{\Needspace{4\baselineskip}\@startsection{subsection}{2}{0pt}{9pt plus 1pt minus 1pt}{3pt}{\sffamily\bfseries\normalsize\color{ink}}}
\makeatother
\setcounter{secnumdepth}{0}
\begin{document}
{\sffamily\small\color{muted}Nonnegative and positive
factorizations\par}
\vspace{3pt}
{\raggedright\hyphenpenalty=10000\exhyphenpenalty=10000\sffamily\bfseries\fontsize{21}{24}\selectfont\color{ink}Rational
factors for rational completely positive boundary matrices\par}
\vspace{7pt}
{\sffamily\small\textbf{Difficulty:} extreme\quad\textbf{Importance:} interesting
to the community\par}
{\sffamily\small\bfseries\color{accent}Status: Lean verified\par}
\vspace{4pt}
{\color{accent}\hrule height 0.8pt}
\vspace{6pt}
\textbf{Rating rationale:} Historical ratings for the original target.
Extreme because rational certificates on arbitrary boundary faces
require new arithmetic control beyond interior geometry; community
importance is exact certification in completely positive optimization.

\textbf{Area:} exact certificates for nonnegative symmetric
factorization

\subsection{Resolution: negative, 13 September
2026}\label{resolution-negative-13-september-2026}

Sidney Holden (Center for Computational Biology, Flatiron Institute,
Simons Foundation) gives a counterexample in
\href{https://github.com/ajt60gaibb/OpenProblemsInNLA/blob/main/references/holden-pf03-2026-09-13/proof/PF03_counterexample.pdf}{Theorem
1.1, proved in Sections 2-6}. The exact integer matrix has order 444,
rank and real cp-rank seven, and strictly positive entries. It lies on
the completely positive boundary and has \textbf{no rational nonnegative
factor of any finite width}. Thus it refutes the universal statement
below; no classification of each smaller order or claim of minimality is
needed.

The matrix is defined by the supplied signed integer matrix \(R\) as
\(A=RR^{\mathsf T}\). An algebraic orthogonal matrix \(O\) gives the
nonnegative real factor \(RO\). Lemma 2.1 excludes all rational
nonnegative factors using a trace-zero quadratic form whose only nonzero
zeros in the certified rational cone lie on seven irrational rays.
Singularity is allowed by the original target and establishes boundary
membership.

A separate
\href{https://github.com/ajt60gaibb/OpenProblemsInNLA/blob/main/references/holden-pf03-2026-09-13/independent-review.md}{independent
informal Codex AI-agent audit} passed the full mathematical argument,
exact facet enumeration and matrix checks. That informal audit supported
\textbf{Solved} under the repository policy and did not claim Lean
verification or external human peer review. The separate formal
verification is recorded below.
\href{https://github.com/ajt60gaibb/OpenProblemsInNLA/blob/main/references/holden-pf03-2026-09-13/README.md}{Author,
verified affiliation, exact certificates and reproduction record} ·
\href{https://github.com/ajt60gaibb/OpenProblemsInNLA/blob/main/references/holden-pf03-2026-09-13/proof/PF03_counterexample.tex}{Proof
source}.

\subsection{Lean proof and verification
evidence}\label{lean-proof-and-verification-evidence}

\textbf{Lean verified --- 19 September 2026 (UTC). Formalization: George
Stepaniants}, Department of Computing and Mathematical Sciences,
California Institute of Technology, with Codex assistance. Sidney Holden
retains the original mathematical and exact seed-data authorship.

The
\href{https://github.com/sgstepaniants/OpenProblemsInNLA/tree/9625a76780183040186664100e24e0e90d8fcc7d/nonnegative-and-positive-factorizations/PF-03/lean}{immutable
Lean proof} proves the full original universal assertion false: some
rational symmetric completely positive boundary matrix of order at least
five has no rational nonnegative factor of \textbf{any positive finite
width}. Boundary is taken in the space of real symmetric matrices. A
proved rational halfspace representation and five zero rows replace the
explicit facet enumeration. This formal witness need not have order 444,
strictly positive entries, or certified minimal cp-rank; those
additional source claims are outside the formalization's scope.

All 25 frozen contracts passed the local serial Lean check and two
separate published-source Linux checks: the
\href{https://github.com/sgstepaniants/OpenProblemsInNLA/actions/runs/35424055075}{fork
run} and
\href{https://github.com/ajt60gaibb/OpenProblemsInNLA/actions/runs/35424087832}{upstream
PR run}. The real Comparator, default kernel, sandbox and rejection
controls passed with only the three standard foundational axioms.
Kernel-mode LeanCert supplies the fixed endpoint certificates used by
the final proof. Two independent nonauthor AI-agent proof reviews and a
separate operational evidence review passed; no official Tau Ceti
service or external human peer review is claimed.

\href{https://github.com/ajt60gaibb/OpenProblemsInNLA/blob/main/nonnegative-and-positive-factorizations/PF-03/lean/README.md}{Formalization,
reproduction and scope} ·
\href{https://github.com/ajt60gaibb/OpenProblemsInNLA/blob/main/nonnegative-and-positive-factorizations/PF-03/lean/verification/linux-2026-09-19/SUMMARY.json}{Exact
execution evidence} ·
\href{https://github.com/ajt60gaibb/OpenProblemsInNLA/blob/main/nonnegative-and-positive-factorizations/PF-03/lean/formalization.yaml}{Formalization
metadata} ·
\href{https://github.com/ajt60gaibb/OpenProblemsInNLA/pull/303}{Upstream
PR}.

\newpage

\subsection{Context and notation}\label{context-and-notation}

A symmetric matrix \(A\) is \textbf{completely positive} if
\(A=BB^\mathsf T\) for an entrywise nonnegative real matrix \(B\) with
finitely many columns. Let \(\mathcal{CP}_n\) be the cone of these
matrices of order \(n\). Its \textbf{cp-rank},
\(\mathop{\mathrm{cpr}}\nolimits(A)\), is the minimum number of columns
in such a factor, with \(\mathop{\mathrm{cpr}}\nolimits(0)=0\).
Boundaries and interiors use the usual Euclidean topology on the vector
space of real symmetric matrices.

\subsection{Problem statement}\label{problem-statement}

For every integer \(n\ge5\) and every

\[A\in\mathbb Q^{n\times n}\cap\partial\mathcal{CP}_n,\]

does there exist a finite integer \(m\ge1\) and a matrix
\(B\in\mathbb Q_{\ge0}^{n\times m}\) such that

\[A=BB^\mathsf T?\]

The width \(m\) is unrestricted. In particular, it need not equal the
real cp-rank or the order of \(A\). This is an existence question for an
exact rational certificate of complete positivity.

Rational matrices in the interior of \(\mathcal{CP}_n\) have rational
completely positive factors. Several boundary classes are also known,
including matrices of rank at most two. The question concerns the
remaining boundary matrices.

\subsection{References}\label{references}

Abraham Berman and Naomi Shaked-Monderer,
\href{https://cot.mathres.org/issues/COT202523.pdf}{\emph{Completely
Positive Matrices over Sets}}, Communications in Optimization Theory
(2025), article 23, §2, Question 2.1 and the explicit boundary Open
Problem, p.~2. Mathieu Dutour Sikirić, Achill Schürmann, and Frank
Vallentin,
\href{https://doi.org/10.1016/j.laa.2017.02.017}{\emph{Rational
factorizations of completely positive matrices}}, Linear Algebra and its
Applications \textbf{523} (2017), 46--51, Theorem 1.1. Max Pfeffer and
José Alejandro Samper,
\href{https://link.springer.com/article/10.1007/s00454-023-00620-y}{\emph{The
Cone of \(5\times5\) Completely Positive Matrices}}, Discrete \&
Computational Geometry \textbf{71} (2024), 442--466, §6, Problem 6.3,
discusses the order-five boundary.

\subsection{Historical status check ---
2026-09-10}\label{historical-status-check-2026-09-10}

Rechecked
\href{https://cot.mathres.org/issues/COT202523.pdf}{Berman--Shaked-Monderer,
§2} and \href{https://arxiv.org/html/2602.05841v2}{Oertel--Schürmann v2,
§6.1.1}, and searched for a 2026 boundary resolution. Rank-at-most-two
and other specified boundary classes have rational factors, but the 2025
paper still poses the general boundary problem. The 2026 interior
theorem does not resolve it. No full boundary construction or
counterexample was located; unrestricted factor width is essential.
\end{document}
